\documentclass[10pt, twocolumn]{article}

\usepackage{amsmath, amssymb}
\usepackage{geometry}
\usepackage{graphicx}
\usepackage{hyperref}
\usepackage{booktabs}
\usepackage{array}
\usepackage{microtype}
\usepackage{setspace}
\usepackage{xcolor}
\usepackage{enumitem}
\usepackage{float}
\usepackage{titlesec}
\usepackage{caption}

\geometry{margin=0.75in, top=0.75in, bottom=0.75in}
\setstretch{1.02}

\titlespacing*{\section}{0pt}{8pt}{4pt}
\titlespacing*{\subsection}{0pt}{6pt}{2pt}
\captionsetup{font=small, labelfont=bf, skip=3pt}

\hypersetup{
    colorlinks=true,
    linkcolor=black,
    urlcolor=blue!70!black,
    citecolor=black
}

% compact lists
\setlist{nosep, leftmargin=1.5em}

% ── title block ──────────────────────────────────────────────────────
\title{\vspace{-1.5em}\textbf{BUICU: Bayesian Belief Updating for ICU\\Crowding Under Uncertainty}\vspace{-0.5em}}
\author{Diego Sanchez\\
\small CS\,109 Challenge Project, Stanford University\\[-2pt]
\small \url{https://buicu-cs109.streamlit.app} \;\;$\mid$\;\;
\url{https://github.com/diegosanchez/buicu}}
\date{}

\begin{document}
\maketitle
\thispagestyle{empty}

% ── Abstract ─────────────────────────────────────────────────────────
\begin{abstract}
\noindent I present BUICU, a fully Bayesian forecasting system for
intensive care unit crowding. Using conjugate Gamma--Poisson inference,
BUICU maintains a continuously updated belief distribution over the
patient arrival rate~$\lambda$ and propagates posterior uncertainty
through 3,000-trajectory Monte Carlo simulation to produce calibrated
occupancy forecasts. On 180 days of synthetic MIMIC-IV--calibrated data
featuring two surge windows, a windowed Bayesian model outscores the
stationary model by 39.4 log-predictive units, a gap entirely
attributable to regime changes. A Law-of-Total-Variance decomposition
reveals that after sufficient observation, 99.4\% of forecast variance
is irreducible Poisson noise, while sensitivity analysis demonstrates
that bed capacity (not arrival-rate estimation) is the dominant driver
of crowding risk. I deploy the model as an interactive Streamlit
application and systematically audit five failure modes, yielding a
combined $\times 2.27$ confidence-interval widening factor that
stretches the crowding probability range from a point estimate of
41.5\% to an honest interval of 18--94\%.
\end{abstract}

% ── 1. Introduction ──────────────────────────────────────────────────
\section{Introduction}

Intensive care units operate at the margin: a single surge can push a
50-bed unit from comfortable to catastrophic overnight, yet most census
projections remain deterministic, projecting false confidence when
uncertainty matters most. BUICU (Belief Updating for ICU Crowding Under
Uncertainty) addresses three questions: (1)~What is $\lambda$, and how
confident is that estimate? (2)~What is the probability of exceeding
capacity within 48~hours? (3)~Which assumptions matter most, and where
does the model break? The answers are exact conjugate inference, Monte
Carlo simulation, and a systematic failure-mode audit, all demonstrated
via an interactive web application.

% ── 2. Methods ───────────────────────────────────────────────────────
\section{Methods}

\subsection{Data Generation}

All experiments use a synthetic dataset generated by a configurable
MIMIC-IV--calibrated simulator. Arrivals follow a \emph{non-homogeneous
Poisson process} with diurnal and weekly modulation, produced via the
\textbf{Lewis--Shedler thinning algorithm}. Two surge windows
(days 30--50 and 110--130, rate $1.8\times$ baseline) introduce
non-stationarity. Length of stay is drawn from a two-component
log-normal mixture (70\% short acute, 30\% long chronic), calibrated to
published MIMIC-IV statistics (median $\approx 2$ days,
mean $\approx 4$ days). Approximately 2\% of discharge timestamps are
randomly censored to simulate realistic data-quality issues. The
generator is fully parameterized: arrival rate, surge windows, capacity,
LOS distribution, and random seed are all configurable, so every result
can be reproduced under alternative scenarios without modifying the
inference pipeline.\footnote{Synthetic data is used for ethical reasons:
real ICU records contain protected health information.}

\subsection{Probabilistic Model}

Let $k_t$ denote the number of admissions on day~$t$. I model each
count as conditionally Poisson:
\begin{equation}\label{eq:likelihood}
    N_t \mid \lambda \;\sim\; \mathrm{Poisson}(\lambda),
\end{equation}
an assumption grounded in the independence and memorylessness of
individual patient deterioration events.\footnote{This independence
assumption is violated during infection waves and trauma surges;
Section~4.1 quantifies the impact.} I place a conjugate
\textbf{Gamma prior} on the unknown rate:
\begin{equation}\label{eq:prior}
    \lambda \;\sim\; \mathrm{Gamma}(\alpha_0,\, \beta_0),
\end{equation}
with $\alpha_0 = 2$, $\beta_0 = 0.2$, giving a prior mean of~10
admissions/day and a standard deviation of~$\approx 7.1$, deliberately
wide to let the data speak. The Gamma--Poisson conjugacy yields a
closed-form posterior after $T$ days of observation:
\begin{equation}\label{eq:posterior}
    \lambda \mid k_{1:T} \;\sim\; \mathrm{Gamma}\!\left(
        \alpha_0 + \textstyle\sum_{t=1}^T k_t,\;\;
        \beta_0 + T
    \right).
\end{equation}
Each morning the update reduces to $\alpha \leftarrow \alpha + k_t$,
$\beta \leftarrow \beta + 1$. No Markov chain Monte Carlo, no
variational approximation is needed; the update is algebraically exact.
The posterior mean $\alpha/\beta$ drifts toward the empirical mean while
the variance $\alpha/\beta^2$ shrinks monotonically, providing a
built-in measure of estimation confidence.

\subsection{Posterior Predictive Distribution}

Integrating out $\lambda$ yields the posterior predictive:
\begin{equation}\label{eq:predpred}
    P(N_{\text{fut}} = k) = \int_0^\infty \!\mathrm{Pois}(k;\lambda)
    \, \mathrm{Gamma}(\lambda;\alpha,\beta)\, d\lambda,
\end{equation}
which evaluates to a \textbf{Negative Binomial} with $r = \alpha$,
$p = \beta/(\beta + \Delta t)$. The NegBin is overdispersed relative
to a Poisson at the posterior mean, reflecting \emph{both} irreducible
Poisson noise \emph{and} residual parameter uncertainty. After
180~days, the 95\% predictive interval is $[8,\, 16]$.

\subsection{Windowed Model for Non-Stationarity}

The stationary model treats all 180~days of history equally. This is
convenient but dangerous: a two-week surge that ended a month ago still
pulls the posterior toward the elevated rate. To let the model
\emph{forget}, I restrict the sufficient statistics to the most recent
$W = 14$ days:
\begin{equation}\label{eq:windowed}
    \alpha_t^{(W)} = \alpha_0 + \!\!\sum_{s=t-W+1}^{t} k_s, \quad
    \beta_t^{(W)} = \beta_0 + W.
\end{equation}
This remains exact conjugate inference over a truncated dataset. The
trade-off is explicit: smaller windows react faster but have higher
variance; larger windows are more stable but slower to detect regime
changes. I adjudicate with the \textbf{one-step-ahead log predictive
score}, a proper scoring rule that cannot be gamed by miscalibrated
predictions:
\begin{equation}\label{eq:logscore}
    S_T = \sum_{t=2}^{T}
    \ln\,\mathrm{NegBin}(k_t;\, \alpha_{t-1},\, p_{t-1}).
\end{equation}

\subsection{Monte Carlo Occupancy Forecasting}

Occupancy $O_t = \sum_i \mathbf{1}[a_i \le t < a_i + L_i]$ depends on
arrival times $a_i$ and lengths of stay $L_i$, so its distribution is
intractable. For each of $S = 3{,}000$ trajectories I:
(i)~sample $\lambda^{(s)}$ from the Gamma posterior;
(ii)~draw arrivals from $\mathrm{Poisson}(\lambda^{(s)} \Delta t)$;
(iii)~bootstrap $L_i^{(s)}$ from the empirical LOS distribution; and
(iv)~overlay the $n_0$ currently observed patients to compute
$O_t^{(s)}$ for $t \in [0,\, 48]$~hours, yielding quantile bands for
operational decision-making.

\subsection{Uncertainty Decomposition}

The \textbf{Law of Total Variance} decomposes forecast uncertainty:
\begin{equation}\label{eq:totvar}
    \mathrm{Var}[N] =
    \underbrace{E[\mathrm{Var}[N \mid \lambda]]}_{\text{aleatoric}}
    + \underbrace{\mathrm{Var}[E[N \mid \lambda]]}_{\text{epistemic}}.
\end{equation}
The aleatoric term ($\Delta t\cdot\alpha/\beta$) is irreducible Poisson
noise. The epistemic term ($\Delta t^2\cdot\alpha/\beta^2$) captures
residual uncertainty about $\lambda$ and vanishes as
$\beta\to\infty$.

% ── 3. Results ───────────────────────────────────────────────────────
\section{Results}

\subsection{Arrival-Rate Estimation and Prior Robustness}

The posterior converges to $\hat\lambda = 11.72$ admissions/day
(95\% credible interval $[11.23,\, 12.23]$). To test robustness, three
priors are run through the same data: \emph{uninformative}
$(\alpha_0, \beta_0) = (0.01, 0.001)$, \emph{weakly informative}
$(2, 0.2)$, and \emph{strong but wrong} $(50, 10)$ (prior mean~5).
After 180~days, pairwise KL divergences between all three posteriors
fall below $10^{-3}$~nats, confirming the Bernstein--von~Mises theorem
in action. The ``wrong'' prior required $\sim$15~days to be overwritten
by the data, a useful calibration of how quickly evidence overwhelms
belief. Figure~\ref{fig:prior_sens} illustrates this convergence.

\begin{figure}[!t]
\centering
\includegraphics[width=\columnwidth]{output/06_prior_sensitivity.png}
\caption{Prior sensitivity: all three priors converge to the same
posterior after $\sim$180 days; the ``wrong'' prior (centered at 5)
is overridden within 15 days.}
\label{fig:prior_sens}
\end{figure}

\subsection{Model Comparison}

The windowed model outscores the stationary model by \textbf{39.4
log-score units}. The gap is entirely attributable to surge periods,
where the stationary model assigns systematically low probability
to elevated counts. Outside surges, the two models perform comparably.
I also compare MLE $\hat\lambda_\text{MLE} = \sum k / T$ against the
full Bayesian posterior. Early in the data (days 1--20), the Bayesian
credible interval is substantially wider than the MLE-based confidence
interval, correctly reflecting the scarcity of evidence; as data
accumulates, the two converge. Posterior predictive $p$-values flag
17 of 180~days as anomalous ($p < 0.025$ or $p > 0.975$); 94\% of
these fall within the surge windows, confirming the model's ability to
detect genuine distributional surprises. Figure~\ref{fig:logscore}
shows the cumulative log-score advantage of the windowed model.

\subsection{Uncertainty Decomposition}

After 180~days the posterior has $\alpha \approx 2100$,
$\beta \approx 180$. The epistemic fraction evaluates to:
\[
    \frac{\Delta t/\beta}{1 + \Delta t/\beta}
    = \frac{1/180}{1 + 1/180}
    \approx 0.6\%.
\]
\textbf{99.4\% of forecast uncertainty is irreducible Poisson noise.}
No amount of additional data can reduce it. The only way to narrow the
predictive interval is to change the underlying process, for example
through triage protocols, regional load-balancing, or elective-surgery
scheduling. This is arguably the most practically important finding:
the bottleneck is not statistical but operational.
Figure~\ref{fig:vardecomp} visualizes how the epistemic component
vanishes as data accumulates.

\subsection{Crowding Forecast and Sensitivity Analysis}

Forecasting from day~36 (a near-capacity snapshot at the onset of the
first surge, with 38 of 50~beds occupied) yields a 48-hour crowding
probability of $\mathbf{41.5\%}$ (bootstrap 95\% CI: 40--43\%), with
peak mean occupancy reaching 44~patients.
Table~\ref{tab:sensitivity} reveals the hierarchy of assumptions.

\begin{table}[H]
\centering
\small
\caption{Sensitivity of 48-hour crowding probability to modeling
assumptions. The dominant lever is bed capacity.}
\label{tab:sensitivity}
\renewcommand{\arraystretch}{1.2}
\begin{tabular}{@{}lcc@{}}
\toprule
\textbf{Scenario} & \textbf{P(crowd.)} & \boldmath$\Delta$\textbf{pp} \\
\midrule
Baseline (stat.\ $\lambda$, emp.\ LOS, 50 beds)
                          & 41.5\% & $-$     \\
Windowed arrival rate     & 52.3\% & $+$10.8  \\
LOS truncated at p90      & 34.1\% & $-$7.4   \\
LOS tails inflated $\times 1.5$ & 48.9\% & $+$7.4   \\
Capacity $-20\%$ (40 beds)& 96.1\% & $+$54.6  \\
Capacity $+20\%$ (60 beds)& 9.0\%  & $-$32.5  \\
\bottomrule
\end{tabular}
\end{table}

\noindent The dominant driver is capacity. A 20\% reduction in available
beds (plausible during a staffing shortage) raises crowding risk from
41\% to 96\%. LOS tail assumptions matter but are secondary; the choice
of arrival-rate model has an intermediate effect. This hierarchy tells
administrators where to focus: \emph{protecting bed availability matters
more than refining the statistical model}.
Figure~\ref{fig:forecast} shows the 48-hour occupancy forecast fan
chart from the day-36 snapshot.

\begin{figure}[!t]
\centering
\includegraphics[width=\columnwidth]{output/04_occupancy_forecast.png}
\caption{48-hour Monte Carlo occupancy forecast (3,000 trajectories)
from day 36. Shaded bands show 50\% and 95\% credible intervals;
the red line marks the 50-bed capacity threshold.}
\label{fig:forecast}
\end{figure}

% ── 4. Discussion ────────────────────────────────────────────────────
\section{Discussion}

\subsection{Failure-Mode Audit}

A model that does not know how it fails is dangerous. I
systematically audit five failure modes, each following the chain
\emph{assumption $\to$ violation $\to$ consequence $\to$ mitigation}:

\smallskip\noindent\textbf{FM1: Non-stationarity.} The Poisson model
assumes constant~$\lambda$; real surge/normal ratio is~1.76. The
windowed model mitigates this (CI~$\times 1.23$).

\smallskip\noindent\textbf{FM2: Independence violations.} Lag-1
autocorrelation $r = 0.08$; LOS--occupancy coupling ignored. A Hawkes
process would address this (CI~$\times 1.02$).

\smallskip\noindent\textbf{FM3: Informative censoring.} 2.1\% missing
discharge timestamps bias LOS downward. Kaplan--Meier correction
required (CI~$\times 1.04$).

\smallskip\noindent\textbf{FM4: Heavy tails.} LOS excess kurtosis of
5.8 means rare 20$+$\,day stays are undersampled by bootstrap.
Generalized Pareto tail fit needed (CI~$\times 1.14$).

\smallskip\noindent\textbf{FM5: Goodhart's Law.} If staff divert
patients based on forecasts, the model learns the \emph{managed}
rate and underpredicts risk when interventions are removed, creating a
dangerous oscillation. Cannot be empirically detected; flagged
unconditionally (CI~$\times 1.05$).

\smallskip\noindent The five penalties compound multiplicatively to a
\textbf{combined widening factor of $\times 2.27$}, stretching the
crowding probability interval to 18--94\%, a stark reminder that a
41\% point estimate carries substantial structural uncertainty beyond
sampling variability.

\subsection{Limitations and Future Directions}

The discharge model ignores the \emph{boarding effect}
($L_i$ independent of occupancy). A Hawkes process would capture
temporal clustering, and a changepoint prior could replace the fixed
window. Validation on real de-identified MIMIC-IV data is needed
before clinical deployment.

\subsection{Ethical Considerations}

Three safeguards are non-negotiable: only synthetic data is used;
every forecast is a distribution, never a point prediction; and failure
modes are front-page content on the website.

\subsection{Use of Large Language Models}

I used an LLM as a \textbf{code reviewer/debugger} (catching
inconsistencies, verifying formulas) and a \textbf{UI/UX collaborator}
(Streamlit CSS, layout, animations). All statistical modeling, Bayesian
inference, mathematical derivations, and analysis were my own work.

% ── 5. Conclusion ───────────────────────────────────────────────────
\section{Conclusion}

BUICU shows that conjugate Bayesian inference with Monte Carlo
simulation produces crowding forecasts that are computationally trivial
and epistemically honest. The central finding, that \emph{bed capacity,
not model refinement, dominates crowding risk}, reframes the problem
from statistical to operational. The deployment at
\url{https://buicu-cs109.streamlit.app} lets visitors manipulate
priors, watch beliefs update, and stress-test capacity scenarios.

% ═══════════════════════════════════════════════════════════════════
% APPENDIX — starts on a new page after the 3-page main body
% ═══════════════════════════════════════════════════════════════════
\clearpage
\onecolumn
\appendix
\renewcommand{\thesection}{\Alph{section}}

\section{Mathematical Derivations}\label{app:math}

\subsection{Conjugate Update}

Combining the Poisson likelihood $P(k \mid \lambda) = \lambda^k e^{-\lambda}/k!$ with the Gamma prior $p(\lambda) \propto \lambda^{\alpha_0 - 1} e^{-\beta_0 \lambda}$ yields
\[
    p(\lambda \mid k_{1:T}) \;\propto\;
    \lambda^{\alpha_0 + \sum k_t - 1}\,
    e^{-(\beta_0 + T)\lambda}
    \;=\;
    \mathrm{Gamma}\!\left(\alpha_0 + \textstyle\sum k_t,\; \beta_0 + T\right).
\]
The posterior mean $\alpha/\beta \to \bar{k}$ as $T \to \infty$, and the posterior variance $\alpha/\beta^2 = O(1/T)$ contracts at the parametric rate.

\subsection{KL Divergence Between Successive Posteriors}

For Gamma posteriors $P_{\text{new}} = \mathrm{Gamma}(\alpha_n, \beta_n)$ and $P_{\text{old}} = \mathrm{Gamma}(\alpha_o, \beta_o)$, the KL divergence admits a closed form:
\begin{align}
    \mathrm{KL}(P_\text{new} \| P_\text{old})
    &= (\alpha_n - \alpha_o)\,\psi(\alpha_n)
    - \ln\Gamma(\alpha_n) + \ln\Gamma(\alpha_o) \nonumber\\
    &\quad + \alpha_o\bigl(\ln\beta_n - \ln\beta_o\bigr)
    + \alpha_n\,\frac{\beta_o - \beta_n}{\beta_n},
\end{align}
where $\psi$ denotes the digamma function. This expression quantifies the information gained per observation.

\subsection{Variance Decomposition}

For $N \mid \lambda \sim \mathrm{Poisson}(\lambda \Delta t)$ with $\lambda \sim \mathrm{Gamma}(\alpha, \beta)$, the Law of Total Variance gives
\begin{align}
    E[\mathrm{Var}[N \mid \lambda]]
    &= \Delta t \cdot \alpha / \beta
    &&\text{(aleatoric),} \\
    \mathrm{Var}[E[N \mid \lambda]]
    &= \Delta t^2 \cdot \alpha / \beta^2
    &&\text{(epistemic).}
\end{align}
The epistemic fraction $(\Delta t / \beta)\big/(1 + \Delta t / \beta) \to 0$ as $\beta \to \infty$, confirming that parameter uncertainty vanishes with additional data while Poisson noise remains irreducible.

\subsection{Random Variables}

Table~\ref{tab:rv} summarizes the random variables and their distributional assumptions.

\begin{table}[H]
\centering
\caption{Random variables used in BUICU.}\label{tab:rv}
\renewcommand{\arraystretch}{1.2}
\small
\begin{tabular}{@{}lll@{}}
\toprule
\textbf{Variable} & \textbf{Distribution} & \textbf{Role} \\
\midrule
$\lambda$ & $\mathrm{Gamma}(\alpha_0, \beta_0)$ & Prior arrival rate \\
$N_t \mid \lambda$ & $\mathrm{Poisson}(\lambda \Delta t)$ & Admission count \\
$\lambda \mid k_{1:T}$ & $\mathrm{Gamma}(\alpha_0\!+\!\Sigma k,\, \beta_0\!+\!T)$ & Posterior rate \\
$N_{\text{fut}}$ & $\mathrm{NegBin}(\alpha,\, \beta/(\beta\!+\!\Delta t))$ & Predictive count \\
$L$ & LogNormal mixture & Length of stay \\
$O_t$ & Monte Carlo & Occupancy \\
\bottomrule
\end{tabular}
\end{table}

% ─────────────────────────────────────────────────────────────────────
\section{Failure-Mode Analysis}\label{app:failure}

Five failure modes are systematically audited, each following the structure: \emph{assumption $\to$ violation $\to$ consequence $\to$ mitigation}. Table~\ref{tab:fm} summarizes the quantitative impact; detailed descriptions follow.

\begin{table}[H]
\centering
\caption{Failure modes and associated CI widening factors.}\label{tab:fm}
\renewcommand{\arraystretch}{1.2}
\small
\begin{tabular}{@{}clc@{}}
\toprule
\textbf{ID} & \textbf{Failure Mode} & \textbf{CI Factor} \\
\midrule
FM1 & Non-stationarity & $\times 1.23$ \\
FM2 & Independence violations & $\times 1.02$ \\
FM3 & Informative censoring & $\times 1.04$ \\
FM4 & Heavy-tailed LOS & $\times 1.14$ \\
FM5 & Goodhart's Law & $\times 1.05$ \\
\midrule
& \textbf{Combined (multiplicative)} & $\boldsymbol{\times 2.27}$ \\
\bottomrule
\end{tabular}
\end{table}

\noindent\textbf{FM1: Non-stationarity.}\;
The constant-$\lambda$ assumption is violated by surge windows in which the arrival rate reaches $1.8\times$ baseline. The stationary posterior averages across regimes, underestimating surge-period risk. A surge/normal ratio of 1.76 exceeds the 1.3 detection threshold. \emph{Mitigation:} windowed posterior (Eq.~\ref{eq:windowed}).

\medskip\noindent\textbf{FM2: Independence violations.}\;
Lag-1 autocorrelation of $r = 0.08$ in daily counts indicates mild temporal clustering. Additionally, LOS may increase during crowding (boarding effect), violating the independence between $L_i$ and $O_t$. \emph{Mitigation:} Hawkes or autoregressive Poisson models.

\medskip\noindent\textbf{FM3: Informative censoring.}\;
Approximately 2.1\% of discharge timestamps are missing. Long-stay patients are disproportionately affected, introducing a downward bias in the empirical LOS distribution. \emph{Mitigation:} Kaplan--Meier survival estimation.

\medskip\noindent\textbf{FM4: Heavy-tailed LOS.}\;
Excess kurtosis of 5.8 implies that rare stays of 20+ days create sustained high-occupancy episodes undersampled by the empirical bootstrap. \emph{Mitigation:} Generalized Pareto tail fit with importance sampling.

\medskip\noindent\textbf{FM5: Goodhart's Law.}\;
If clinical staff use forecasts to divert patients or accelerate discharges, the model learns the \emph{managed} rate and will underpredict risk when interventions are removed. This mode is empirically undetectable and is flagged unconditionally.

\medskip\noindent The combined widening factor of $\times 2.27$ stretches the 48-hour crowding probability from a point estimate of 41.5\% to an adjusted interval of 18--94\%.

% ─────────────────────────────────────────────────────────────────────
\section{CS\,109 Concepts Demonstrated}\label{app:concepts}

Table~\ref{tab:concepts} maps the 16 CS\,109 topics exercised in this project to their specific realizations.

\begin{table}[H]
\centering
\caption{CS\,109 concepts and their application in BUICU.}\label{tab:concepts}
\renewcommand{\arraystretch}{1.15}
\small
\begin{tabular}{@{}p{3.2cm}p{12cm}@{}}
\toprule
\textbf{Category} & \textbf{Concepts Applied} \\
\midrule
Probability Foundations
  & Random variables ($N_t$, $L$, $O_t$, $\lambda$); probability distributions (Poisson, Gamma, Negative Binomial, LogNormal); conditional probability $P(N_t \mid \lambda)$; Bayes' theorem via conjugate updating. \\
\addlinespace
Bayesian Inference
  & Posterior predictive distribution; conjugate priors with closed-form updates; Law of Total Variance decomposition; Monte Carlo simulation (3{,}000 trajectories). \\
\addlinespace
Frequentist Methods
  & MLE $\hat{\lambda} = \Sigma k / T$; CLT-based confidence intervals; KL divergence for information-gain quantification; posterior predictive $p$-value testing. \\
\addlinespace
Model Evaluation
  & Proper scoring rules (log predictive score); calibration via coverage probabilities and PIT histograms; sensitivity analysis; prior robustness analysis. \\
\bottomrule
\end{tabular}
\end{table}

% ─────────────────────────────────────────────────────────────────────
\section{Interactive Application}\label{app:website}

The deployed Streamlit application (\url{https://buicu-cs109.streamlit.app}) is organized as a guided narrative with five modules:

\begin{enumerate}[leftmargin=2em]
    \item \textbf{Arrival-rate elicitation.} Users submit a prior guess via slider; the posterior mean is then revealed with a comparison metric.
    \item \textbf{Prior explorer.} Adjustable $(\alpha_0, \beta_0)$ sliders update the prior--posterior overlay in real time.
    \item \textbf{Sequential update.} A 180-day scrubber animates the posterior evolution, displaying KL divergence and anomaly flags at each step.
    \item \textbf{Crowding forecast.} Capacity and horizon controls trigger cached Monte Carlo recomputation; a sensitivity chart displays $P(\text{crowding})$ at 80\%, 100\%, and 120\% capacity.
    \item \textbf{Model evaluation.} Tabbed panels compare stationary vs.\ windowed log scores, MLE vs.\ Bayesian intervals, and an expandable failure-mode audit with detection evidence.
\end{enumerate}

% ─────────────────────────────────────────────────────────────────────
\section{Supplementary Figures}\label{app:figures}

\begin{figure}[H]
\centering
\includegraphics[width=0.75\textwidth]{output/10_log_score_comparison.png}
\caption{Cumulative one-step-ahead log predictive score. The windowed model outperforms the stationary model by 39.4 units, with the gap concentrated in the surge windows (shaded).}
\label{fig:logscore}
\end{figure}

\begin{figure}[H]
\centering
\includegraphics[width=0.75\textwidth]{output/12_variance_decomposition.png}
\caption{Law of Total Variance decomposition over 180 days. The epistemic component (orange) vanishes as $\lambda$ is learned; the aleatoric floor (blue) is irreducible Poisson noise.}
\label{fig:vardecomp}
\end{figure}

\end{document}
